\section{The acquisition--authority envelope}
\label{sec:envelope}

The acquisition/closure split becomes substantially more useful when it is
lifted from a design rule to a comparison theory.  A scientific discovery
controller has four coupled responsibilities: generate or select a route,
acquire evidence through it, convert acquired material into decision-usable
state while preserving useful future cues, and decide whether the task may
close.  A retriever improves only one part of this object; a stopping policy can
save cost while making another part worse.  This section defines the joint
object and derives the conditions under which such local improvements can be
composed into a scientifically stronger controller.

\subsection{Worlds, obligations and a non-compensatory frontier}

Let $w\in\mathcal W$ be a scientific world and let $h_t$ be the complete
observable transcript after charged step $t$.  The task contract assigns to
$w$ a set $\mathcal O(w)$ of \emph{material scientific obligations}.  An
obligation may be a claim that needs evidence, a contradiction that needs both
sides inspected, a dataset or code artifact needed to establish reproducibility,
a route whose availability must be resolved, or a previously retrieved source
that must be processed under the current question.  Materiality means that
resolving the obligation can change the predeclared decision, conclusion,
uncertainty state or admissible closure state.  It is not synonymous with
relevance to a search query.

A controller $\pi$ comprises route generation $G$, acquisition $A$, a
question-conditioned state compiler $M$, and closure authority $C$.  Under
budget $B$, source-access contract $\mathcal I$ and task distribution $P$, let
its outcome vector be
\[
 Z_B(\pi,w)=(D,U,V,-F,Y,-K).
\]
Here $D$ is decision-relevant evidence discovery, $U$ is final scientific
utility, $V$ is the value of search options preserved for prospectively hidden
future questions, $F$ is false task closure, $Y$ is valid closure yield, and
$K$ is fully charged cost.  Signs are chosen so larger is better except that
$F$ is a non-compensatory safety coordinate: a gain in report quality cannot
purchase permission to hide a false closure.

\begin{definition}[Acquisition--authority envelope]
\label{def:acquisition-authority-envelope}
For an admissible policy class $\Pi$,
\[
 \mathcal E_{B,\mathcal I}(\Pi)
 =\operatorname{cl}\left\{
     \mathbb E_{w\sim P}[Z_B(\pi,w)] :
     \pi\in\Pi,\ \mathbb E[K(\pi,w)]\le B
   \right\}.
\]
Its \emph{guarded frontier} retains only points satisfying the prospectively
fixed false-closure, valid-closure-yield, provider-validity and cost guards.
A candidate establishes envelope improvement only by lying outside the guarded
frontier of the frozen donor-complete class under the same $P$, $B$ and
$\mathcal I$; winning against a conveniently weak member of that class is
insufficient.
\end{definition}

This definition absorbs the strongest neighbouring components instead of
opposing them.  Lexical and corpus-side augmentation, reasoning-aware retrieval,
fielded inspection and selective fetching, and citation expansion enlarge the
action family of $G$ and $A$ \cite{sage2026,agentir2026,sieve2026,rethinkinglitsearch2026}.
Claim--evidence coverage, marginal information utility, conformal continuation,
decision-aware and evidence-aware stopping enlarge the policy and signal class
\cite{halt2026,deepcontrol2026,micp2026,decisionstop2026,donotstopearly2026,confidencebasedstop2026}.
Finite source certificates, obligation graphs and question-conditioned memory
enlarge $M$ and the observable closure state
\cite{knowplan2026,icore2026,memchain2026,scienceintent2026}.  These are donor
territory.  The residual question is whether their composition reaches a better
guarded frontier and whether its closure decision is identifiable from the
state the product actually retains.

\begin{proposition}[Donor-monotonicity]
\label{prop:donor-monotonicity}
If $\Pi_1\subseteq\Pi_2$ under identical task, access and budget contracts, then
$\mathcal E_{B,\mathcal I}(\Pi_1)\subseteq
\mathcal E_{B,\mathcal I}(\Pi_2)$.  Consequently, adding a sound donor action,
signal or memory operation cannot make the attainable envelope smaller, though
it may leave the guarded frontier unchanged after its cost is charged.
\end{proposition}

\begin{proof}
Every policy feasible in $\Pi_1$ remains feasible in $\Pi_2$; take the closure
of their attainable outcome vectors.  The guarded statement follows because
the same non-compensatory filters are applied to the superset.
\end{proof}

This elementary monotonicity result changes the novelty test.  A new paper that
improves retrieval should be absorbed into the comparator class.  The theory is
falsified, not protected, if the expanded donor class reaches every proposed
candidate point under matched resources.

\subsection{The acquisition ceiling: why seven emitted identifiers matter}

Let $G_w$ be a complete evaluator gold set for a fixed task world, $A_\pi(w)$
the union of gold-normalized identities ever emitted by any admissible route
during a frozen run, and $S_\pi(w)$ the final submitted set.  A downstream
ranker, reader or synthesizer that is forbidden to invent evidence satisfies
$S_\pi(w)\cap G_w\subseteq A_\pi(w)\cap G_w$.

\begin{theorem}[Run-conditional acquisition ceiling]
\label{thm:acquisition-ceiling}
For any such downstream policy and any nonempty $G_w$,
\[
 \operatorname{Recall}(S_\pi;G_w)
 \leq \frac{|A_\pi(w)\cap G_w|}{|G_w|}.
\]
More generally, if obligations have nonnegative weights $\mu_o$ and each
obligation can be discharged only by at least one member of evidence set $E_o$,
then the unnormalized resolved-obligation weight is at most the total weight of
obligations for which $A_\pi(w)\cap E_o\ne\varnothing$.
\end{theorem}

\begin{proof}
The submitted relevant set is a subset of the acquired relevant set, so its
cardinality cannot be larger.  For the weighted form, an obligation whose
required evidence set is disjoint from the acquired pool cannot be discharged
by a downstream operation over that pool; summing the remaining nonnegative
weights gives the bound.
\end{proof}

The theorem separates candidate-generation failures from selection failures
without assuming that a particular retriever is optimal.  In the archived
external adverse campaign, all routes emitted only 7 of the 229 gold
identifiers; six were submitted and one was discarded.  The run-conditional
recall ceiling was therefore $7/229=0.030568$, and 222 gold identities were
unavailable to every downstream selection rule on that capture.  This does not
estimate general web discoverability.  It proves the narrower but decisive
causal fact that reranking or synthesis could not repair that run.  The adverse
terminal and the separate provider-validity failure remain unchanged.

The same analysis exposes why matched request counts are not a matched
scientific opportunity.  Let $\xi(a,G_w)$ be the admissible exposure of action
$a$: the prospectively validated opportunity, through the scorer-compatible
interface, to emit identities in $G_w$.  Equal $\sum_a 1$ does not imply equal
$\sum_a\xi(a,G_w)$.

\begin{proposition}[Call matching does not identify exposure matching]
\label{prop:call-exposure}
For every positive call budget $m$, there exist two $m$-call policies under the
same nominal budget whose run-conditional acquisition ceilings are respectively
one and zero.  Hence call-count equality alone gives no nontrivial bound on the
difference in discoverability.
\end{proposition}

\begin{proof}
Take a one-item gold set.  Give the first policy $m$ calls whose admissible
support contains the item with a deterministic first-call return, and give the
second $m$ calls to an interface whose output support is disjoint from the gold
identity namespace.  Both spend $m$ calls; their acquired-gold fractions are one
and zero.
\end{proof}

The archived campaign gave the lexical arm three calls to
the only scoring-admissible backend and the governed arm one such call plus two
inadmissible calls.  Its three-versus-three request count was therefore not an
exposure match.  In a future protected study, exposure must be certified
outcome-blind through interface-class witnesses or treated as
\textsc{cannot-check}; gold itself cannot be used as a worker-visible route
oracle.

\subsection{When can a product implement scientific closure?}

Let $s:\mathcal W\rightarrow\mathcal S$ be everything a donor product makes
available to its closure rule at a fixed decision time: acquired evidence,
coverage and utility signals, route states, compiled memory and charged
resource state.  Write
\[
 \mathcal Y_C=\{\textsc{close},\textsc{open},\textsc{cannot-check}\}
\]
and let $\kappa:\mathcal W\rightarrow\mathcal Y_C$ be the correct task
decision under the declared scientific contract.

\begin{theorem}[Closure factorization]
\label{thm:factorization}
There exists a decision map $g:\mathcal S\rightarrow\mathcal Y_C$ such that
$\kappa=g\circ s$ if and only if $\kappa$ is constant on every fibre
$s^{-1}(z)$.
\end{theorem}

\begin{proof}
If $\kappa=g\circ s$ and $s(w)=s(w')$, then
$\kappa(w)=g(s(w))=g(s(w'))=\kappa(w')$.  Conversely, if $\kappa$ is constant
on every fibre, define $g(z)$ to be $\kappa(w)$ for any $w$ satisfying
$s(w)=z$.  Fibre constancy makes the definition independent of the chosen
$w$.
\end{proof}

The theorem gives an exact necessary-and-sufficient condition, not a claim that
a particular architecture is intrinsically more expressive.  Coverage,
confidence and route-local stop signals are sufficient for task closure only
when no two worlds with the same retained state require different task
decisions.  A state that drops the materiality or unresolved status of an
unavailable route creates a mixed fibre and cannot be repaired by a different
classifier over that same state.  An information-equivalent product that
retains the missing coordinate has pure fibres and can tie.  This is exactly the
boundary exhibited by the protected exact-contract battery: the available-route
product fails on unresolved material routes, while the ideal typed product is
decision-identical to the authority rule.

\subsection{The price of an observationally unresolved world}

The factorization obstruction has a quantitative version.  Consider two task
worlds $w_0,w_1$ whose correct decisions are respectively \textsc{close} and
\textsc{open}.  Let $P_0$ and $P_1$ be the distributions of the complete
observable transcript available to the controller before its decision.  Count
\textsc{cannot-check} as not yielding a valid closure in $w_0$ but not as false
closure in $w_1$.

\begin{theorem}[Indistinguishable-world lower bound]
\label{thm:indistinguishable}
For every randomized closure rule, its false-closure probability $F_1$ in
$w_1$ and valid-closure shortfall $L_0$ in $w_0$ satisfy
\[
 F_1+L_0\geq 1-\operatorname{TV}(P_0,P_1).
\]
In particular, if the two worlds are observationally identical, then
$F_1+L_0=1$ and $\max(F_1,L_0)\geq 1/2$.
\end{theorem}

\begin{proof}
The decision rule is a binary test between $P_0$ and $P_1$ once
\textsc{open} and \textsc{cannot-check} are combined as ``do not close.''  The
sum of its two testing errors is bounded below by one minus total variation.
When $P_0=P_1$, a transcript-dependent close probability has the same
expectation in both worlds, so false closure plus failure to close the closable
world equals one.
\end{proof}

This lower bound prevents a cosmetic conversion of every negative result into a
positive terminal.  If the decisive route is absent from the observable action
support, no stopping heuristic can create the missing distinction.  Scientific
progress must instead do at least one of three things: invent an acquisition
action that makes $P_0$ and $P_1$ distinguishable, retain a state coordinate
that an earlier compiler erased, or introduce and validate a substantive world
assumption that rules out one member of the pair.  Otherwise a calibrated open
or \textsc{cannot-check} state is the only honest output.

\subsection{Authority completion and its maximality}

For a transcript $h$, let $R(h)$ be the registered set of material obligations
that are not yet independently discharged.  Assume an \emph{adversarial
realizability contract}: for each $o\in R(h)$ there exists a task world
consistent with $h$ in which the evidence behind $o$ changes the scientific
decision or admissible closure state.  Define the obligation-completed closure
rule
\[
 C^+(h)=
 \begin{cases}
  \textsc{close}, & R(h)=\varnothing\ \text{and the task contract is valid},\\
  \textsc{cannot-check}, & \text{required observability or provider validity is unestablished},\\
  \textsc{open}, & \text{otherwise}.
 \end{cases}
\]

\begin{theorem}[Maximally permissive sound completion]
\label{thm:maximal-completion}
Relative to the declared obligation and realizability contract, $C^+$ has no
false closure.  Moreover, any rule using the same contract that closes at a
history where $R(h)\ne\varnothing$ is unsound on at least one compatible world.
Thus no universally sound rule closes a strict superset of the histories closed
by $C^+$.
\end{theorem}

\begin{proof}
If $C^+$ closes, every registered material obligation is discharged and the
contract is valid, which is the contract's definition of admissible closure.  If
another rule closes with $o\in R(h)$, adversarial realizability supplies a world
consistent with the same history in which $o$ changes the decision; closure is
false in that world.  Therefore every universally sound rule must refuse that
history.
\end{proof}

The theorem is broad but explicitly contract-relative.  It does not claim that
an open-web system can enumerate every material obligation from finite search.
Completeness of $R(h)$ requires either a closed registry, a validated sampling
model, independently curated hidden-route custody, or a declared
\textsc{cannot-check} boundary.  The theorem nevertheless provides a positive
composition result: any future acquisition, ranking, stopping or memory donor
can be installed below $C^+$ without losing sound closure, and any improvement
that discharges more obligations moves the controller toward valid closure
rather than merely increasing a surrogate score.

\begin{corollary}[When acquisition enlarges closure authority]
Adding a donor mechanism can expand discovery or reduce cost by
Proposition~\ref{prop:donor-monotonicity}.  It enlarges the set of histories that
can be soundly closed only if the resulting observation/state map splits a
previously mixed closure fibre or discharges an obligation in $R(h)$.
\end{corollary}

This is the sense in which the proposed theory envelops its nearest work.
Retrieval, coverage, information utility, stopping, obligation certificates and
memory are all admissible mechanisms inside the envelope.  Their scientific
value is measured by the joint frontier they make reachable; their outputs do
not acquire global authority merely by being numerically confident.

\subsection{A failure-derived candidate: obligation-directed adaptive
exploration}

The theory also specifies a materially stronger acquisition candidate rather
than treating fail-closed authority as a substitute for finding evidence.  Let
$m_t(o)$ be the posterior unresolved mass of material obligation $o$, let
$Q_t(a)$ be the prospectively evaluated option value that action $a$ preserves
for hidden future questions, and let $J_t(a)$ be the expected reduction in
closure-fibre ambiguity---the probability mass of closable/open worlds that the
action is expected to separate.  An \emph{obligation-directed adaptive
exploration} (ODAE) policy chooses from route invention, retrieval, processing,
contradiction resolution, state preservation/recovery and closure verification
by a charged index of the form
\[
 a_t\in\arg\max_{a\in\mathcal A(h_t)}
 \frac{
   \operatorname{LCB}_t\!\left[
     \Delta\sum_{o\in R(h_t)}\mu_o m_t(o)\mid a
   \right]
   +\lambda Q_t(a)+\beta J_t(a)
 }{\operatorname{cost}_t(a)}.
\]
The lower-confidence term prevents an uncertain route estimate from being
treated as established gain; $Q_t$ prevents current-question compression from
destroying future route cues; and $J_t$ values actions that make the closure
decision identifiable even when their immediate document count is low.  If no
action has validated support through the frozen interface, the controller emits
\textsc{cannot-check} rather than assigning that route zero value.

Generic value-of-information, active search, coverage estimation and adaptive
stopping are donor-owned.  The proposed residual is their non-compensatory
coupling to material scientific obligations, option preservation and closure
fibre separation under one charged budget.  ODAE is a prospective policy family,
not an executed result or current general capability.  Its direct falsifier is a
donor-complete controller that reaches the same or a better guarded frontier on
fresh tasks.  Its causal tests must separately vary (i) route support, (ii)
retrieved-but-unprocessed evidence, (iii) contradiction sides, (iv) future-query
optionality and (v) closable clean cases; improvement on only the family that
inspired the policy is insufficient.

\subsection{A prospective cross-arena discriminator}

The theory converts the historical adverse terminals into explicit research
problems rather than deleting them.  The 7/229 exposure result demands route
invention and admissible-support estimation before reranking.  The one-third
scoring-eligible exposure mismatch demands source-interface matching rather
than call-count matching.  The retrieved-but-unprocessed and future-query
failures demand option-valued state preservation.  Provider-invalid and
unresolved-route failures demand independent observability custody and an
authority-complete state.  Clean closable tasks remain a required counterweight
against an always-open policy.

A separately frozen successor gives these problems a common empirical target.
It specifies 768 source-disjoint task-world clusters---192 in each of four
strata---across four arenas, identical admissible route exposure, three frozen
primary comparators, an intersection requirement over every contrast, and a
joint power target of 0.90 under its declared planning model.  That design is a
protocol, not a result.  Archived baseline-task-world reproduction, four
source-disjoint arena snapshots, a capability-matched external comparator,
independent evaluator/gold custody, and live response/cost receipts are still
required.  Until those bindings exist and the protected campaign runs, the
general cross-arena envelope-superiority hypothesis remains untested and the
historical manuscript terminal remains unchanged.
